\section{Every fixed width at one hidden layer}\label{sec:fixed-width}
The affine reduction above does not extend to several hidden neurons. Nevertheless, exact input-permutation symmetry combined with a higher-order approximation gives a subexponential result at every fixed hidden width. In this section there is one hidden layer of width $n$, and
\[
 z(y)=b^\top\tanh(Ay),\qquad y\in\{-1,1\}^{p},\quad p=d+1.
\]
The original cube is the section $y_0=1$. The canonical Gaussian energy and variances are
\[
 \mathcal E_n(A,b)=\frac dn\|A\|_F^2+n\|b\|^2,
 \quad A_{ji}\sim N(0,n\sigma^2/d),
 \quad b_j\sim N(0,\sigma^2/n).
\]
As before, condition first on absolute output margin greater than $1/2$ everywhere and then on exact training labels. Write $E_n(g)$ for the energy infimum of a function event.

\begin{theorem}[Fixed-width shallow networks]\label{thm:fixed-width}
Fix $n\ge1$. Suppose the target has a depth-two, width-one implementation in which every weight has absolute value at most one and the signed output margin is greater than $1/2$ everywhere. At learner width $n$, the strong-prior posterior has, on every nonanticipating input sequence,
\begin{equation}\label{eq:fixed-width-budget}
 \sum_tq_t\le
 \exp\!\left(O_n\!\left(d^{\frac n{n+2}}\log^4(d+2)\right)\right).
\end{equation}
The constant may depend on the fixed learner width $n$. The same bound, plus one, holds for all sufficiently small positive scales. Expected posterior-sampling mistakes, majority mistakes, confidence failures, and high-probability sampled mistakes obey the consequences of Theorem~\ref{thm:main}.

More generally, for fixed $n$ and fixed integer $q\ge1$, a target energy bound $E_n(f)\le K$ gives a strong-prior budget with logarithm at most
\begin{equation}\label{eq:general-width-budget}
 O_{n,q}\!\left(
 \left[1+\left(\frac{K^2}{d}+K^{1+1/(q+1)}\right)^{\frac n{n+2}}\right]
 \log(d+K+2)\right).
\end{equation}
For the targets of the first assertion, whose energy is $K=O(d)$,
each fixed $\epsilon>0$ and a sufficiently large fixed Taylor order give
the alternative exponent
\[
 O_{n,\epsilon}(d^{n/(n+2)+\epsilon}\log(d+2)).
\]
\end{theorem}

The fixed-width hypothesis matters. This is not a bound uniform over widths that grow with $d$. For widths $n=2$ and $n=3$, the exponents of $d$ in~\eqref{eq:fixed-width-budget} are respectively $1/2$ and $3/5$, up to the displayed logarithms. For $n=1$, the scalar argument already gives the sharper logarithmic factor in Corollary~\ref{cor:one}.

\subsection{Close input columns force an exact Boolean symmetry}
Every function with $E_n(g)\le K$ has a weak-margin representative of that energy. Indeed, the energy is coercive and its weak-margin cell is closed; increasing the readout slightly shows that weak and strict infima agree. Define
\[
 B=\sqrt{K/n},\qquad E_0=nK/d.
\]
The representative satisfies $\|b\|\le B$ and $\sum_i\|A_i\|^2\le E_0$, where $A_i$ denotes its $i$th column. The map $s\mapsto b^\top\tanh(s)$ is $B$-Lipschitz.

Choose a radius $\rho\le1/(8B)$, select a maximal $\rho$-separated subset of the columns, and assign every column to a center within distance $\rho$. Thus each group lies in a radius-$\rho$ ball. Swapping two coordinates from the same group changes the preactivation by at most $4\rho$ and the output by at most $4B\rho\le1/2$. Two opposite labels with absolute margin at least $1/2$ would require an output difference at least one. Hence the Boolean function is invariant under this swap, and therefore under every permutation within each group.

Let there be $r$ groups. The centers are original columns, so their squared norms sum to at most $E_0$. At least $r/2$ lie in the radius-$\sqrt{2E_0/r}$ ball. Packing radius-$\rho/2$ balls about these centers gives
\[
 \frac r2\le\left(1+\frac{2\sqrt{2E_0/r}}{\rho}\right)^n.
\]
If $r>2^{n+1}$, subtracting one from the $n$th root of the left side loses at most a factor two. Rearranging, and including the complementary case, gives
\begin{equation}\label{eq:column-group-count}
 r\le C_n\left[1+(E_0\rho^{-2})^{n/(n+2)}\right].
\end{equation}

% Independent lemma for fixed-width orbit compression.
\begin{lemma}[Permutation moments and Taylor approximation]\label{lem:perm-moments}
Let $A\in\mathbb R^{n\times p}$ have its columns partitioned into groups,
with every group contained in a Euclidean ball of radius $\rho$.
For a fixed $y\in\{-1,1\}^p$, independently permute its coordinates
within each group to obtain $Y$. Write $S=AY$, $\mu=\mathbb ES$, and
$U=S-\mu$. Then, for every row $j$, every real $t$, and every integer $s\ge1$,
\begin{align}
 \mathbb E e^{tU_j}&\le \exp(8p\rho^2t^2),\label{eq:perm-mgf}\\
 \mathbb E|U_j|^s&\le(8\sqrt{s}\sqrt p\,\rho)^s.\label{eq:perm-moment}
\end{align}
Consequently, if $\|b\|_2\le B$ and $q\ge1$ is an integer, the function
\[
 P(y)=\sum_{j=1}^n b_j
       \sum_{a=0}^{q}\frac{\tanh^{(a)}(\mu_j)}{a!}
                         \mathbb E U_j^a
\]
satisfies
\begin{equation}\label{eq:perm-taylor}
 \left|\mathbb E\bigl[b^T\tanh(AY)\bigr]-P(y)\right|
 \le \sqrt n B\bigl(16\sqrt{q+1}\sqrt p\,\rho\bigr)^{q+1}.
\end{equation}
\end{lemma}
\begin{proof}
Within a group of size $m$, center the scalar entries in row $j$ as
$z_1,\ldots,z_m$, so $\sum_i z_i=0$. Their range has length at most
$2\rho$, since the original group lies in a radius-$\rho$ ball.
The centered contribution to $U_j$ is twice the sum of a uniformly
selected $k$-subset of these entries, where $k$ is the number of positive
coordinates of $y$ in the group. Its moment generating function is
\[
 \frac{e_k(e^{2tz_1},\ldots,e^{2tz_m})}{\binom mk}
 \le\left(\frac1m\sum_{i=1}^m e^{2tz_i}\right)^k.
\]
Here $e_k$ is the elementary symmetric polynomial. The inequality follows
by repeatedly averaging two of its nonnegative arguments while preserving
their sum: the coefficient of their product is nonnegative, so this
operation cannot decrease $e_k$, and its maximum at a fixed sum occurs
when all arguments are equal. The cases $k=0,m$ cause no difficulty.

A uniform draw $Z$ from the centered entries has mean zero and lies in an
interval of length at most $2\rho$. The elementary bounded-variable
Hoeffding lemma gives $\mathbb E e^{2tZ}\le\exp(2\rho^2t^2)$; even the
looser range bound $4\rho$ gives $\exp(8\rho^2t^2)$, which we use here.
Thus the group contribution has moment generating function at most
$\exp(8m\rho^2t^2)$. Independence across groups proves
\eqref{eq:perm-mgf}.

Chernoff's argument gives
$\mathbb P(|U_j|>u)\le2\exp(-u^2/(32p\rho^2))$.
For an integer $\ell\ge1$, integration yields
\[
 \mathbb E|U_j|^{2\ell}\le2(32p\rho^2)^\ell\ell!
 \le2(32\ell p\rho^2)^\ell.
\]
Take $\ell=\lceil s/2\rceil\le s$ and use monotonicity of $L^r$ norms.
The resulting bound is at most $(8\sqrt s\sqrt p\rho)^s$, proving
\eqref{eq:perm-moment}.

For completeness, $|\tanh^{(s)}(u)|\le2^s s!$ for all $u\in\mathbb R$:
write the $s$th derivative as $P_s(\tanh u)$, start from $P_0(v)=v$,
and use $P_{s+1}(v)=(1-v^2)P_s'(v)$. Its coefficient $\ell^1$ norm
is at most $2^s s!$ by induction, since $\deg P_s=s+1$.
Taylor's theorem therefore bounds the order-$q$ remainder by
$2^{q+1}|U_j|^{q+1}$. Apply \eqref{eq:perm-moment}, sum with
weights $|b_j|$, and use $\|b\|_1\le\sqrt n B$ to obtain
\eqref{eq:perm-taylor}.
\end{proof}

\paragraph{Choice of radius.}
Set $s=q+1$ and
\[
 \rho=\min\left\{\frac1{8B},\;
      \frac{(8\sqrt n B)^{-1/s}}{16\sqrt{s}\sqrt p}\right\}.
\]
Then swaps within a group change the original output by at most
$4B\rho\le1/2$, giving exact Boolean invariance under all such swaps.
The Taylor error in \eqref{eq:perm-taylor} is at most $1/8$.
If the original output has signed margin at least $1/2$ on the full cube,
then the permutation average has the same margin, and the approximating
function $P$ has signed margin at least $3/8$ on every input.
Moreover
\[
 \rho^{-2}\le C_{n,q}\left(B^2+pB^{2/(q+1)}\right),
\]
where $C_{n,q}$ is a constant depending only on the two fixed integers.

\subsection{A finite description of the symmetrized Taylor polynomial}
Fix the multiplicities $m_1,\ldots,m_r$ of the column groups. For row $j$ and group $g$, record the power sums
\[
 P_{g,j,\ell}=\sum_{i\in G_g} A_{ji}^{\ell},\qquad 1\le\ell\le q.
\]
Together with the $n$ readout weights, these are $nqr+n$ real descriptors. Their number is linear in $r$ when $n,q$ are fixed.

We verify that they determine every moment used in the Taylor approximation. For one group, write its scalar entries as $a_1,\ldots,a_m$ and suppose the queried sign vector has $k$ positive coordinates in that group. A uniform permutation is equivalent to choosing a uniform $k$-subset, with indicators $\xi_i$. Put $V=\sum_i a_i\xi_i$ and $P_\ell=\sum_i a_i^\ell$. For $h\le q$,
\begin{equation}\label{eq:subset-moment}
 \E V^h=\sum_{\pi\in\mathcal P([h])}
 \frac{(k)_{|\pi|}}{(m)_{|\pi|}}
 \sum_{\substack{i_B\in[m]\ (B\in\pi)\\\text{all distinct}}}
       \prod_{B\in\pi}a_{i_B}^{|B|}.
\end{equation}
Here $(t)_v$ is a falling factorial, and terms with $|\pi|>m$ are zero. The formula partitions ordered index tuples according to their coincidence pattern. Each coefficient $(k)_v/(m)_v$ lies in $[0,1]$.

If a partition has block sizes $s_1,\ldots,s_v$, its distinct-index sum is a polynomial in $P_1,\ldots,P_h$. Explicitly, inclusion--exclusion on coincidence patterns gives
\begin{equation}\label{eq:distinct-powers}
 \sum_{i_1,\ldots,i_v\text{ distinct}}\prod_{t=1}^v a_{i_t}^{s_t}
 =\sum_{\tau\in\mathcal P([v])}
   \left[\prod_{C\in\tau}(-1)^{|C|-1}(|C|-1)!\right]
   \prod_{C\in\tau}P_{\sum_{t\in C}s_t}.
\end{equation}
This is the usual partition-lattice inclusion--exclusion; it also follows inductively by subtracting all nontrivial coincidence patterns from the unrestricted product of power sums. Thus~\eqref{eq:subset-moment} has degree at most $h$ in the descriptors, and its coefficient magnitudes and number of terms are bounded by constants depending only on $h$.

The group's centered preactivation is
\[
 U_g=2\left(V-\frac{k}{m}P_1\right).
\]
Its moments through order $q$ are likewise polynomials of degree at most their order, with coefficients uniformly bounded in $m,k$. Different groups are permuted independently. Expanding the moments of $U=\sum_gU_g$ therefore gives polynomials of degree at most $h$ in all group descriptors, with at most $C_q r^q$ terms and coefficients bounded by $C_q$, for $h\le q$. The mean is
\[
 \mu_j=\sum_g\left(\frac{2k_g}{m_g}-1\right)P_{g,j,1}.
\]
Consequently the function
\begin{equation}\label{eq:descriptor-predictor}
 T(y)=\sum_{j=1}^n b_j
       \sum_{h=0}^q\frac{\tanh^{(h)}(\mu_j)}{h!}\E U_j^h
\end{equation}
is completely determined by the group sizes, these descriptors, and the number $k_g$ of positive signs in each queried group. It is not required to be a neural network: its role is to encode the Boolean function, which is already determined by its uniformly separated sign.

\subsection{Uniform quantization of the descriptors}
Let $A_* =\sqrt{E_0}$ and
\[
 H=p\max\{1,A_*\}^q.
\]
Every power-sum descriptor has absolute value at most $H$, and every readout weight has absolute value at most $B$. Quantize all of them on a common grid with mesh $\eta\le1$. We give the elementary uniform continuity estimate needed for this step, including its dependence on the growing dimension $p$.

For fixed $n,q$, the displayed moment polynomials have degree at most $q$, at most $C_qp^q$ monomials, and coefficients bounded by $C_q$. On the descriptor box of radius $H+1$, changing each coordinate by at most $\eta$ changes any such polynomial by at most
\[
 C_q p^q (H+1)^q\eta.
\]
This follows by telescoping the factors of each monomial; the constants absorb its degree and number of terms. The mean $\mu_j$ changes by at most $p\eta$. Every derivative of tanh of any fixed order is bounded on the real line: inductively it is a polynomial in tanh, using $(\tanh)'=1-\tanh^2$. Applying the mean value theorem to the coefficients $\tanh^{(h)}(\mu_j)$ and then changing the readout weights therefore gives
\begin{equation}\label{eq:descriptor-lipschitz}
 |T(y)-\widetilde T(y)|
 \le C_{n,q}(1+B)p^{q+1}(H+1)^q\eta
\end{equation}
for every full-cube input $y$, with one constant independent of all group sizes and sign counts. This deliberately loose bound also controls the magnitudes of the moment polynomials appearing when the readout is changed.

Choose
\[
 \eta=\min\left\{1,
 [16C_{n,q}(1+B)p^{q+1}(H+1)^q]^{-1}\right\}.
\]
The output perturbation is at most $1/16$. If the Taylor approximation has signed margin at least $3/8$, its quantized version retains margin at least $5/16$ and exactly the same Boolean function.

There are at most
\[
 M=3+\left\lceil\frac{2\max\{H,B\}}{\eta}\right\rceil
\]
grid values per descriptor. Since $E_0=nK/d$, $B=\sqrt{K/n}$ and $p=d+1$, for fixed $n,q$ this obeys
\[
 \log M=O_{n,q}(\log(d+K+2)).
\]
One can put every group consecutively by an input permutation. The group sizes have at most $p^r$ possibilities, and the quantized descriptors have at most $M^{nqr+n}$ possibilities. Choose one grouping and one quantized description deterministically for each function orbit. A shared description would define the same canonical Boolean function by~\eqref{eq:descriptor-predictor}; hence two distinct orbits cannot share it. The number of exact signed-augmented-coordinate symmetry orbits below energy $K$ is therefore at most
\[
 J_{n,q}(d,K)\le
 \sum_{r=1}^{\min\{p,\lceil R_{n,q}\rceil\}}
      p^r M^{nqr+n},
\]
where the group count $R_{n,q}$ comes from the energy packing estimate. This count permits invalid descriptor lists and repeated encodings, both of which only overcount.


\subsection{Orbit count, Gaussian posterior, and the growing Taylor order}
The radius above and~\eqref{eq:column-group-count} imply
\[
 r\le C_{n,q}\left[1+
 \left(\frac{K^2}{d}+K^{1+1/(q+1)}\right)^{n/(n+2)}\right].
\]
Here $p/d\le2$ has been absorbed into the constant. The finite descriptor count therefore gives
\[
 \log J_{n,q}(d,K)
 \le O_{n,q}\left(
 \left[1+\left(\frac{K^2}{d}+K^{1+1/(q+1)}\right)^{n/(n+2)}\right]
 \log(d+K+2)\right).
\]
Every exact signed-augmented-coordinate orbit has at most $2^pp!$ functions, and all its functions have equal Gaussian mass at every positive scale. The proof is the same parameter isometry used in the scalar case, applied simultaneously to the first-layer columns. Gaussian small-noise rates give a posterior tail at most $\beta$ outside $E_n(g)\le E_n(f)$ after every target-consistent history. Lemma~\ref{lem:blocks} consequently yields
\[
 \sum_tq_t\le J_{n,q}(d,K)[1+\log(2^pp!)]+\beta2^d.
\]
The polynomial orbit-size factor is absorbed in~\eqref{eq:general-width-budget}; the tail vanishes in a strong-prior limit. This proves the fixed-$q$ assertion with all posterior prediction criteria.

For clarity we now track the constants needed when the Taylor order grows. The number of partitions of a set of size $q$ is at most $q^q$, and every factorial occurring in~\eqref{eq:distinct-powers} is at most $q!$. Expanding~\eqref{eq:subset-moment}, the centering binomial, and the sum over independent groups therefore bounds the relevant coefficients and the constants in~\eqref{eq:descriptor-lipschitz} by
\[
 \exp(O(q\log(q+1)))
\]
apart from a factor depending on the fixed $n$. One way to see that combining groups does not produce an extra exponential in $q^2$ is to expand the $h$ ordered factors directly: there are at most $r^h$ assignments of their group indices, and the product of their groupwise partition counts and factorial bounds is at most $\exp(O(h\log(h+1)))$, since the group degrees sum to $h\le q$. The tanh derivative bound in Lemma~\ref{lem:perm-moments} has the same factorial size.

Such a target, with every weight bounded by one, is an augmented signed majority with affine optimum $a_*$ satisfying $\|a_*\|^2\le p$. Set $a=a_*/\sqrt d$ and $b_*=1/[2\tanh(1/\sqrt d)]$. Since $\tanh(t)\ge t/2$ for $0\le t\le1$, its weak-margin scalar energy is
\[
 d\|a\|^2+b_*^2\le p+d=2d+1.
\]
Cloning it $n$ times with readouts $b_*/n$ preserves $\mathcal E_n$ exactly. Thus we may take $K=2d+1$, so $E_0=O_n(1)$ and $B=O_n(\sqrt p)$.

Choose $q=\lceil\log(p+1)\rceil$ and $s=q+1$. The explicit radius gives
\[
 \rho^{-2}
 =\max\{64B^2,\,256s p(8\sqrt n B)^{2/s}\}
 =O_n(p\log(p+1)).
\]
Consequently
\[
 r=O_n\!\left((p\log(p+1))^{n/(n+2)}\right).
\]
Since $H=p\max\{1,\sqrt{E_0}\}^q$, we have $\log(H+1)=O_n(\log(p+1))$. The explicit quantization mesh and the coefficient bounds just established give
\[
 \log M=O_n(\log^2(p+1)).
\]
There are $nqr+n=O_n(r\log(p+1))$ descriptors and at most $p^r$ group-size lists. Hence
\[
 \log J_{n,q}(d,2d+1)
 =O_n\!\left(p^{n/(n+2)}
                \log^{3+n/(n+2)}(p+1)\right).
\]
This implies the slightly simpler bound~\eqref{eq:fixed-width-budget}, completing the theorem. No growing-order approximation is needed for the separately stated fixed-$q$ conclusion.
